\subsection{Best-effort Observability}

Deep learning has gained prominence in tandem with growth in training set sizes, parameter counts, and training FLOPs \citep{marcus2018deep,kaplan2020scaling}.
Since, AI/ML workloads have brought to market powerful next-generation compute accelerators affording up to hundreds of thousands of processing cores (e.g., by Tenstorrent, Groq, Graphcore, SambaNova, and Cerebras) \citep{zhang2016cambricon,emani2021accelerating,jia2019dissecting,medina2020habana}.
These accelerator devices are anticipated to drive advances in both agent-based modeling (ABM) and high-performance computing (HPC) writ large \citep{perumalla2022computer,vanessendelft2025record}, and opportunity exists to benefit digital evolution.

To maximize compute throughput, emerging accelerator platforms exploit highly-distributed dataflow architectures --- though this can come at the cost of on-device memory scarcity and locality.
For instance, to pack 880,000 processors onto a single chip, the Cerebras Wafer-Scale Engine (WSE) platform provisions only 48kB memory per processor, introduces multi-hop message routing, and allows data export only from the chip periphery.
These constraints mirror broader trends in HPC where data storage and movement are becoming more expensive relative to compute \citep{buluc2021randomized,khan2021analysis,gholami2024ai}.

In some regards, computational costs and constraints arising around data storage and movement resemble the logistical costs and constraints faced in research studying real-world systems.
Rather than ``complete,'' deterministic data observability, most real-world studies instead rely on dynamic, query-driven sampling --- a paradigm we term ``inferential observability.''
Such approaches can degrade gracefully under resource scarcity and intermittent disruptions.

In scoping allowed relaxations, distinction can be drawn between \textit{core computation} and \textit{observation}.
In scenarios where core computation has iterative structure, runtime irregularities compound and --- in the worst case --- can produce a pathological cascade.
Although systematic biases may still be introduced, irregularity in observation affects only \textit{post hoc} analysis, where it may often be more readily reasoned about and accounted for.
Possible benefits of best-effort observation strategies, by contrast, include efficiencies in dynamically exploiting convenient data locality or available openings in unused bandwidth and memory, as well as fault robustness and capped-budget utilization guarantees.
Such an \textit{inferential observability} paradigm thus strikes an intermediate compromise, which may suit a broader swath of use cases.

Effectively harnessing AI/ML accelerators for scientific computing thus poses substantial, but broadly pertinent, engineering challenges.
In the case of digital evolution, a key challenge is tracking population dynamics across a vast, highly-distributed fabric of memory-constrained processors.
Such data collection can be essential to digital evolution work.
Likewise, in the context of application-oriented evolutionary computation (EC), diagnostic information is necessary to troubleshoot, tune parameters, and develop theory \citep{hernandez2023suite,hooker1995testing}.

It has repeatedly been the case that significant steps forward in biological knowledge piggyback on innovations making new types or quantities of data available, technological or otherwise.
For instance, evolutionary theory arose in the context of burgeoning taxonomic collections \citep{winsor2009taxonomy}, core ideas in pathology and developmental biology arose from microscopy \citep{turner1890cell}, high-throughput sequencing has made symbiosis a key concept in organismal biology (e.g., gut microbiome) \citep{durack2019gut}, and new imaging technologies are driving new questions about ecological interactions across continental-scale distance \citep{stark2016toward}.
Although factors besides feasibility have also catalyzed significant advances in life science (e.g., Mendel, Redi, Semmelweis), growth in capabilities to generate and collect data play a longstanding and ongoing role in enabling new biological inquiry \citep{strasser2012data}.
Contemporary biology enjoys profound, ongoing gains in data availability \citep{sulston1983embryonic,sheth2017multiplex,weeks2023deep}, nevertheless fundamental limitations exist in data completeness, particularly with respect to historical accounts of natural history \citep{benton2011assessing,delsuc2005phylogenomics}.

In contrast to biology, digital evolution has, from the outset, enjoyed perfect fidelity and absolute completeness in data collection.
Indeed, such observability is a key driver of scientific interest in using \textit{in silico} models for evolution research \citep{o2003digital}.

Although some domains of genetic programming have been highlighted for their capability to produce intuitive symbolic expressions \citep{hu2023genetic,javed2022simplification}, it is also the case that discerning the functionality of some evolved artifacts can require extensive experiment-driven analyses.
Such knockout trials --- in contrast to other aspects of digital evolution --- are notable in that combinatoric tractability has held back completely exhaustive analysis for all but the smallest genomes \citep{nitash2021information}.
Knockout assay experiments have, therefore, typically limited to single-site \citep{adami2006digital}, all-pairs \citep{kumawat2023fluctuating}, or iterative approaches \citep{langdon2014improving,moreno2021case}.

Massively parallel and distributed processing power, which has been argued crucial to future directions in digital evolution \citep{moreno2022best,taylor2016open}, peels away the tractability of digital evolution's existing perfect observability paradigm.
One concern is that many-process experiments can produce greater volumes of data than is feasible to store, much less analyze \citep{klasky2021data}.
For instance, even under serial processing, maintaining full records of genetic program instruction history under sexual recombination has proven to be a highly technically demanding, data-intensive task \citep{mcphee2016using}.
Parallel and distributed computing also introduces challenges in runtime overhead from communication and synchronization required for data collection and introduces the possibility of data loss when components fail \citep{snir2014addressing}.
Continuing with the phylogenetic example, typical tracking approaches are sensitive to even small amounts of data loss and, in a distributed computing context, require runtime inter-process communication to reclaim memory from extinct lineages \citep{moreno2024algorithms}.

Fortunately for digital evolution, research in biology, by necessity, already routinely works around issues of incomplete and imperfect data.
As such, existing methods can provide a valuable foothold in scenarios where combinatorial effects or runtime multiprocessing make exhaustive direct observation impractical.
This section reviews work leveraging methods borrowed from biology to negotiate data limitations on both fronts: 1) application of mark-recapture approaches from ecology to characterize fitness landscapes and 2) application of reconstruction-based approaches inspired by bioinformatics for robust, decentralized phylogeny tracing.

In WSE-based experiments, agent migrations across the chip pose substantial logistical challenges in tracking evolutionary history as it unfolds \citep{moreno2024analysis}.
On WSE, instrumentation must be lightweight because memory used comes at the direct cost of simulation capacity.
Even were sufficient memory and bandwidth available, a complete evolutionary history --- upwards of a quadrillion replication events in some simulations \citep{singhvi2025scalable} --- would be of negligible tractability in its entirety; instead, means are needed to extract a representative summary.

In evolutionary studies, one key aspect of history is phylogeny --- the structure of ancestry relationships among organisms \citep{faith1992conservation, stamatakis2005phylogenetics,french2023host,kim2006discovery,lenski2003evolutionary}.
In addition to tracing the history of evolutionary events \citep{dawkins2016ancestor,hedges2015tree,hinchliff2015synthesis}, phylogenetic analysis can test more general questions about the underlying mode and tempo of evolution \citep{felsenstein1985phylogenies}.
Classically, these analyses have been applied to species-level macroevolution; however, population- and organism-level dynamics can also be inferred \citep{genthon2023cell, levy2015quantitative,stadler2013recovering,Moreno2025ecology,lewinsohn2023statedependent,nozoe2017inferring}.
Phylogeny data can also be leveraged in proactive efforts to influence the trajectory of an evolving population \citep{scott2018inferring}.
In application-oriented EC, phylogenetic information can help guide evolution toward desired outcomes \cite{lalejini2024phylogeny,lalejini2024runtime,murphy2008simple,burke2003increased}; it also can provide informative diagnostics to practitioners \cite{hernandez2022can,shahbandegan2022untangling}.
Finally, digital phylogenies additionally serve as a useful testbed for bioinformatics because of capability to configure ground truth conditions directly \cite{daudey2024aevol,haller2023slim,moreno2025extending}.

To collect ancestry data from WSE experiments, we have adopted strategies inspired by real-world experiments that assume data collection is expensive, limited, and non-deterministic.
Rather than exhaustively tracking parent-child relationships \cite{dolson2023phylotrackpy}, we adopt an approach inspired by bioengineered lineage tracking \cite{moreno2022hereditary} --- in which progressive DNA barcode inserts identify closely-related organisms by shared commonalities \cite{masuyama2019dna}.
We use analogous barcoding with single-bit markers, budgeting a 64-bit capacity per genome \cite{moreno2025testing}.
Once budgeted capacity fills, barcodes overwrite each other to maintain an approximate record \cite{moreno2024algorithms}.
% Barcode origin times can be inferred from storage position, allowing direct estimation of lineage divergence
% Barcodes remain chronologically attributed, allowing estimation of the point in time when lineages diverged from conflicts \cite{}.
Phylogenies incorporating multiple organisms can be constructed by agglomerative trie-building over shared barcode prefixes \cite{moreno2023toward}.

\citet{moreno2027trust}
Developments in high-performance computing (HPC) technology continue to drastically increase quantities of available processor cores, driven in large part by compute-dense hardware accelerator devices.
In the context of digital evolution, this explosive growth offers opportunity to advance both hypothesis-driven explorations of multi-scale biological phenomena and application-driven evolutionary optimization targeting hard problem domains.
Available processing power, however, is by nature vastly distributed --- a challenge exemplified by emerging next-generation AI/ML hardware accelerator platforms, such as the 880,000-processor Cerebras Wafer-Scale Engine (WSE).
Strict constraints in on-device data storage and movement, compounded by vulnerability to component failures, may favor partial relaxations of the traditional deterministic computing paradigm.
Such best-effort relaxations, however, complicate reproducibility and risk introducing artifactual biases.
We explore these concerns, developing a framework to measure runtime behavior of best-effort code and examining case studies of best-effort computing in digital evolution projects.
The first case study applies best-effort CPU-cluster multiprocessing to a digital multicell workload, which provides 92\% scaling efficiency at 64 processes ($2.1\times$ speedup) and exhibits robust median quality of service, even under hardware anomalies.
The second case study examines WSE-based simulations, demonstrating best-effort strategies to track spatiotemporal population history --- through sparse, asynchronous device-to-host sampling that tolerates hardware faults.
In sum, across potential forms and scopes of best-effort relaxation, we argue that digital evolution is uniquely positioned to contribute in developing post-deterministic HPC paradigms.
